\documentclass[UTF8,fontset=fandol,11pt]{ctexart}
\usepackage[a4paper,margin=24mm]{geometry}
\usepackage{amsmath,amssymb,amsthm,graphicx,tabularx,booktabs,hyperref}
\newtheorem{definition}{定义}[section]
\newtheorem{theorem}[definition]{定理}
\newtheorem{example}[definition]{例题}
\newtheorem{exercise}[definition]{练习}
\hypersetup{unicode=true,colorlinks=true,linkcolor=blue}
\title{函数极限与连续：本科一年级微积分讲义计划}
\date{}
\begin{document}
\maketitle
\section{一元函数极限：定义、左右极限与无穷远极限}\label{sec:S01}
% UNIT S01.U01
\begin{definition}
设函数 $f$ 在点 $x_0$ 的某个去心邻域内有定义。若存在常数 $A$，使得对任意 $\varepsilon>0$，都存在 $\delta>0$，满足当 $0<|x-x_0|<\delta$ 时有 $|f(x)-A|<\varepsilon$，则称 $A$ 为 $f(x)$ 当 $x\to x_0$ 时的极限，记为
\[
\lim_{x\to x_0}f(x)=A.
\]
其中，$x_0$ 不要求属于函数的定义域。\label{def:point-limit}
\end{definition}
% UNIT S01.U02
\begin{definition}
设函数 $f$ 在 $x_0$ 的某个左侧去心邻域内有定义。若存在常数 $A_-$，使得对任意 $\varepsilon>0$，都存在 $\delta>0$，当 $x_0-\delta<x<x_0$ 时有 $|f(x)-A_-|<\varepsilon$，则称 $A_-$ 为 $f(x)$ 当 $x\to x_0^-$ 时的左极限，记为
\[
\lim_{x\to x_0^-}f(x)=A_-.
\]
设函数 $f$ 在 $x_0$ 的某个右侧去心邻域内有定义。若存在常数 $A_+$，使得对任意 $\varepsilon>0$，都存在 $\delta>0$，当 $x_0<x<x_0+\delta$ 时有 $|f(x)-A_+|<\varepsilon$，则称 $A_+$ 为 $f(x)$ 当 $x\to x_0^+$ 时的右极限，记为
\[
\lim_{x\to x_0^+}f(x)=A_+.
\]
\label{def:one-sided-limit}
\end{definition}
% UNIT S01.U03
\begin{definition}
设函数 $f$ 在某个右半无限区间 $(M,+\infty)$ 上有定义。若存在常数 $A$，使得对任意 $\varepsilon>0$，都存在 $R>M$，当 $x>R$ 时有 $|f(x)-A|<\varepsilon$，则称 $A$ 为 $f(x)$ 当 $x\to+\infty$ 时的极限，记为
\[
\lim_{x\to+\infty}f(x)=A.
\]
类似地，若函数在某个左半无限区间 $(-\infty,M)$ 上有定义，且对任意 $\varepsilon>0$ 都存在 $R<M$，使得当 $x<R$ 时有 $|f(x)-A|<\varepsilon$，则记为
\[
\lim_{x\to-\infty}f(x)=A.
\]
\label{def:infinite-limit}
\end{definition}
% UNIT S01.U04
\begin{theorem}
设函数 $f$ 在点 $x_0$ 的某个去心邻域内有定义，则
\[
\lim_{x\to x_0}f(x)=A
\]
当且仅当
\[
\lim_{x\to x_0^-}f(x)=A
\quad\text{且}\quad
\lim_{x\to x_0^+}f(x)=A.
\]
也就是说，函数在一点的双侧极限存在，当且仅当该点的左极限和右极限都存在且相等。\label{thm:two-sided-limit-criterion}
\end{theorem}
% UNIT S01.U05
\begin{proof}
先设 $\lim_{x\to x_0}f(x)=A$。给定任意 $\varepsilon>0$，由双侧极限的定义，存在 $\delta>0$，使得当 $0<|x-x_0|<\delta$ 时有 $|f(x)-A|<\varepsilon$。特别地，当 $x_0-\delta<x<x_0$ 或 $x_0<x<x_0+\delta$ 时，上述不等式分别成立，因此左极限和右极限都等于 $A$。

反过来，设左极限和右极限都等于 $A$。给定任意 $\varepsilon>0$，存在 $\delta_1>0$，使得当 $x_0-\delta_1<x<x_0$ 时有 $|f(x)-A|<\varepsilon$；存在 $\delta_2>0$，使得当 $x_0<x<x_0+\delta_2$ 时有 $|f(x)-A|<\varepsilon$。取 $\delta=\min\{\delta_1,\delta_2\}$，则当 $0<|x-x_0|<\delta$ 时，$x$ 必在 $x_0$ 的左侧或右侧，从而都有 $|f(x)-A|<\varepsilon$。因此 $\lim_{x\to x_0}f(x)=A$。
\end{proof}
% UNIT S01.U06
\begin{theorem}
设函数 $f$ 在点 $x_0$ 的某个去心邻域内有定义。若
\[
\lim_{x\to x_0}f(x)=A
\quad\text{且}\quad
\lim_{x\to x_0}f(x)=B,
\]
则 $A=B$。因此，函数在同一点处的有限极限至多有一个。\label{thm:limit-uniqueness}
\end{theorem}
% UNIT S01.U07
\begin{proof}
假设 $A\ne B$。取 $\varepsilon=|A-B|/3>0$。由 $\lim_{x\to x_0}f(x)=A$，存在 $\delta_1>0$，使得当 $0<|x-x_0|<\delta_1$ 时有 $|f(x)-A|<\varepsilon$；由 $\lim_{x\to x_0}f(x)=B$，存在 $\delta_2>0$，使得当 $0<|x-x_0|<\delta_2$ 时有 $|f(x)-B|<\varepsilon$。

取 $\delta=\min\{\delta_1,\delta_2\}$，并取一个满足 $0<|x-x_0|<\delta$ 的 $x$。由三角不等式，
\[
|A-B|\leq |A-f(x)|+|f(x)-B|<2\varepsilon=\frac{2}{3}|A-B|,
\]
这与 $|A-B|>0$ 矛盾。因此 $A=B$。
\end{proof}
% UNIT S01.U08
\begin{example}
计算
\[
\lim_{x\to2}(3x-1)\text{。}
\]
为说明双侧极限的处理方法，分别考察 $x\to2^-$ 和 $x\to2^+$ 两种情形。
\end{example}
% UNIT S01.U09
\par\noindent\textbf{解：}当 $x\to2^-$ 或 $x\to2^+$ 时，都有 $3x-1\to3\cdot2-1=5$。更直接地，对任意 $\varepsilon>0$，取 $\delta=\varepsilon/3$。当 $0<|x-2|<\delta$ 时，
\[
|(3x-1)-5|=|3x-6|=3|x-2|<3\delta=\varepsilon\text{。}
\]
因此由极限定义，
\[
\lim_{x\to2}(3x-1)=5\text{。}
\]\par
% UNIT S01.U10
\begin{example}
设
\[
f(x)=
\begin{cases}
1,&x<0,\\
2,&x>0.
\end{cases}
\]
判断 $\lim_{x\to0}f(x)$ 是否存在。
\end{example}
% UNIT S01.U11
\par\noindent\textbf{解：}当 $x\to0^-$ 时，始终有 $f(x)=1$，所以
\[
\lim_{x\to0^-}f(x)=1\text{。}
\]
当 $x\to0^+$ 时，始终有 $f(x)=2$，所以
\[
\lim_{x\to0^+}f(x)=2\text{。}
\]
由于左右极限不相等，根据定理 \ref{thm:two-sided-limit-criterion}，$\lim_{x\to0}f(x)$ 不存在。\par
% UNIT S01.U12
\begin{exercise}
\label{mathagent:exercise:bf60311c135cf0c14f3cb42069a931ed175758be7d89dff861bed83a8d332803}
设
\[
f(x)=
\begin{cases}
2x+1,&x<1,\\
5,&x>1.
\end{cases}
\]
判断 $\lim_{x\to1}f(x)$ 是否存在；若存在，求出该极限。
\end{exercise}
\section{极限的基本性质与运算}\label{sec:S02}
% UNIT S02.U01
\begin{definition}
设函数 $f,g,h$ 在点 $x_0$ 的某个去心邻域内有定义。若在该邻域内有 $h(x)\leq f(x)\leq g(x)$，且存在有限数 $L$，使得
\[
\lim_{x\to x_0}h(x)=L,\qquad \lim_{x\to x_0}g(x)=L,
\]
则
\[
\lim_{x\to x_0}f(x)=L\text{。}
\]
这称为夹逼原理。\label{def:squeeze-principle}
\end{definition}
% UNIT S02.U02
\begin{theorem}
设函数 $f$ 与 $g$ 在点 $x_0$ 的某个去心邻域内有定义，且 $\lim_{x\to x_0}f(x)=A$，$\lim_{x\to x_0}g(x)=B$。则
\[
\lim_{x\to x_0}(f(x)+g(x))=A+B,
\qquad
\lim_{x\to x_0}(f(x)-g(x))=A-B,
\]
\[
\lim_{x\to x_0}f(x)g(x)=AB.
\]
若 $B\neq0$，并且在某个去心邻域内 $g(x)\neq0$，则
\[
\lim_{x\to x_0}\frac{f(x)}{g(x)}=\frac{A}{B}.
\]
\label{thm:limit-algebra}

复合极限的基础结论：设 $g$ 在 $x_0$ 的某个去心邻域内有定义，$\lim_{x\to x_0}g(x)=b$；外函数 $f$ 在 $b$ 的某个去心邻域内有定义，且 $\lim_{t\to b}f(t)=L$。若在 $x_0$ 的某个去心邻域内，$g(x)\neq b$ 且 $g(x)$ 属于 $f$ 的定义域（因而复合函数有定义），则
\[
\lim_{x\to x_0}f(g(x))=L.
\]
这里要求内函数最终不取到 $b$，不能仅凭上述两个极限就无条件推出复合极限。
\end{theorem}
% UNIT S02.U03
\begin{proof}
先证明和的法则。给定 $\varepsilon>0$，由 $f(x)\to A$ 与 $g(x)\to B$，分别存在 $\delta_1>0$ 与 $\delta_2>0$，使得当 $0<|x-x_0|<\delta_1$ 时 $|f(x)-A|<\varepsilon/2$，当 $0<|x-x_0|<\delta_2$ 时 $|g(x)-B|<\varepsilon/2$。取 $\delta=\min\{\delta_1,\delta_2\}$，则
\[
\left|(f(x)+g(x))-(A+B)\right|\leq |f(x)-A|+|g(x)-B|<\varepsilon.
\]
故 $f(x)+g(x)\to A+B$。由于 $|(-g(x))-(-B)|=|g(x)-B|$，由定义有 $-g(x)\to-B$；差的法则由 $f(x)-g(x)=f(x)+(-g(x))$ 及和的法则得到。

先说明局部有界性：若 $f$ 在 $x_0$ 的某个去心邻域内有定义，且有有限极限 $A$，则 $f$ 在某个去心邻域内有界。事实上，在极限定义中取误差为 $1$，可得 $\delta_0>0$，使得当 $0<|x-x_0|<\delta_0$ 时，
\[
|f(x)-A|<1,\qquad |f(x)|\leq |A|+|f(x)-A|<|A|+1.
\]
这只涉及去心邻域，不要求 $f(x_0)$ 有定义。

对于积，令 $M=|A|+1>0$。给定 $\varepsilon>0$，由两个极限的定义可分别取 $\delta_1,\delta_2>0$，使得当 $0<|x-x_0|<\delta_1$ 时 $|g(x)-B|<\varepsilon/(2M)$，当 $0<|x-x_0|<\delta_2$ 时 $|f(x)-A|<\varepsilon/(2(|B|+1))$。取 $\delta=\min\{\delta_0,\delta_1,\delta_2\}$，则当 $0<|x-x_0|<\delta$ 时，
\[
|f(x)g(x)-AB|\leq |f(x)|\,|g(x)-B|+|B|\,|f(x)-A|
<\frac{\varepsilon}{2}+\frac{|B|\varepsilon}{2(|B|+1)}<\varepsilon.
\]
故由极限定义，$f(x)g(x)\to AB$。

若 $B\neq0$，由 $g(x)\to B$，可取 $\delta_3>0$，使得当 $0<|x-x_0|<\delta_3$ 时 $|g(x)-B|<|B|/2$。于是
\[
|g(x)|\geq |B|-|g(x)-B|>\frac{|B|}{2}>0,
\]
故分母不仅非零，而且远离零，并有
\[
\left|\frac{1}{g(x)}-\frac{1}{B}\right|=\frac{|g(x)-B|}{|g(x)||B|}\leq\frac{2|g(x)-B|}{|B|^2}.
\]
给定 $\varepsilon>0$，再由极限定义取 $\delta_4>0$，使得当 $0<|x-x_0|<\delta_4$ 时 $|g(x)-B|<\varepsilon|B|^2/2$。当 $0<|x-x_0|<\min\{\delta_3,\delta_4\}$ 时，上述倒数误差小于 $\varepsilon$。所以 $1/g(x)\to1/B$。再由已证明的积的法则，得到
\[
\frac{f(x)}{g(x)}=f(x)\frac{1}{g(x)}\to\frac{A}{B}.
\]

最后证明复合极限结论。按该结论的假设，取 $r>0$，使得当 $0<|x-x_0|<r$ 时，$g(x)\neq b$ 且 $f(g(x))$ 有定义。给定 $\varepsilon>0$，由 $\lim_{t\to b}f(t)=L$，存在 $\eta>0$，使得当 $0<|t-b|<\eta$ 时 $|f(t)-L|<\varepsilon$。由 $g(x)\to b$，存在 $\delta_5>0$，使得当 $0<|x-x_0|<\delta_5$ 时 $|g(x)-b|<\eta$。因此，当 $0<|x-x_0|<\min\{r,\delta_5\}$ 时，
\[
0<|g(x)-b|<\eta,\qquad |f(g(x))-L|<\varepsilon.
\]
由极限定义即得 $\lim_{x\to x_0}f(g(x))=L$。
\end{proof}
% UNIT S02.U04
\begin{theorem}
设函数 $f,g,h$ 在点 $x_0$ 的某个去心邻域内有定义，并且在该邻域内满足
\[
h(x)\leq f(x)\leq g(x).
\]
若
\[
\lim_{x\to x_0}h(x)=L,
\qquad
\lim_{x\to x_0}g(x)=L,
\]
则
\[
\lim_{x\to x_0}f(x)=L.
\]
\label{thm:squeeze-theorem}
\end{theorem}
% UNIT S02.U05
\begin{proof}
给定 $\varepsilon>0$。由 $h(x)\to L$，存在 $\delta_1>0$，使得当 $0<|x-x_0|<\delta_1$ 时 $h(x)>L-\varepsilon$；由 $g(x)\to L$，存在 $\delta_2>0$，使得当 $0<|x-x_0|<\delta_2$ 时 $g(x)<L+\varepsilon$。取 $\delta=\min\{\delta_1,\delta_2\}$，则当 $0<|x-x_0|<\delta$ 时，
\[
L-\varepsilon<h(x)\leq f(x)\leq g(x)<L+\varepsilon.
\]
因此 $|f(x)-L|<\varepsilon$，从而 $\lim_{x\to x_0}f(x)=L$。
\end{proof}
% UNIT S02.U06
\begin{theorem}
设函数 $f$ 在点 $x_0$ 的某个去心邻域内有定义，且 $\lim_{x\to x_0}f(x)=A$。若 $A>0$，则存在 $\delta>0$，使得当 $0<|x-x_0|<\delta$ 时 $f(x)>0$；若 $A<0$，则存在 $\delta>0$，使得当 $0<|x-x_0|<\delta$ 时 $f(x)<0$。\label{thm:limit-sign-preservation}
\end{theorem}
% UNIT S02.U07
\begin{proof}
设 $A>0$。取 $\varepsilon=A/2$。由极限定义，存在 $\delta>0$，使得当 $0<|x-x_0|<\delta$ 时 $|f(x)-A|<A/2$，于是
\[
f(x)>A-\frac{A}{2}=\frac{A}{2}>0.
\]
当 $A<0$ 时，取 $\varepsilon=|A|/2$，则
\[
f(x)<A+\frac{|A|}{2}=\frac{A}{2}<0.
\]

\end{proof}
% UNIT S02.U08
\begin{example}
计算
\[
\lim_{x\to1}\frac{x^2-1}{x-1}.
\]
要求在保持定义域条件 $x\neq1$ 的前提下进行代数变形。
\end{example}
% UNIT S02.U09
\par\noindent\textbf{解：}当 $x\neq1$ 时，$x^2-1=(x-1)(x+1)$，所以
\[
\frac{x^2-1}{x-1}=x+1.
\]
原函数虽然在 $x=1$ 处没有定义，但极限只考察 $x\neq1$ 且 $x\to1$ 时的函数值。因此
\[
\lim_{x\to1}\frac{x^2-1}{x-1}=\lim_{x\to1}(x+1)=1+1=2.
\]\par
% UNIT S02.U10
\begin{example}
设 $f(x)=x^2-3x+2$。求 $\lim_{x\to2}f(x)$，并说明所得极限的符号。
\end{example}
% UNIT S02.U11
\par\noindent\textbf{解：}由极限的和、差与积法则，
\[
\lim_{x\to2}(x^2-3x+2)=2^2-3\cdot2+2=0.
\]
因此该极限等于 $0$。由于极限值为零，保号定理不能据此断言 $f(x)$ 在 $2$ 的去心邻域内恒正或恒负。\par
% UNIT S02.U12
\begin{exercise}
\label{mathagent:exercise:b4c9c9a1fb82bcef052fdbf6e1539831ed88cee016dbcbf83a50e5e8d2311be0}
设函数 $f$ 与 $g$ 在点 $x_0$ 的某个去心邻域内有定义，且
\[
\lim_{x\to x_0}f(x)=2,
\qquad
\lim_{x\to x_0}g(x)=-1.
\]
求
\[
\lim_{x\to x_0}\frac{f(x)+g(x)}{f(x)-g(x)}.
\]
并说明使用商的极限法则时需要检查的条件。
\end{exercise}
\section{无穷小、基本极限与等价无穷小}\label{sec:S03}
% UNIT S03.U01
\begin{definition}
设函数 $\alpha(x)$ 在考察点 $x_0$ 的某个去心邻域内有定义。若
\[
\lim_{x\to x_0}\alpha(x)=0,
\]
则称 $\alpha(x)$ 是当 $x\to x_0$ 时的无穷小。这里的“无穷小”表示函数值趋近于零，而不表示函数在某个邻域内恒不为零。若对任意 $M>0$，都存在 $\delta>0$，使得当 $0<|x-x_0|<\delta$ 时有
\[
|\alpha(x)|>M,
\]
则称 $\alpha(x)$ 是当 $x\to x_0$ 时绝对值意义下的无穷大，即 $|\alpha(x)|\to+\infty$。本节的“无穷大”采用这一含义，它本身不指定函数值的正负方向。

正无穷极限与负无穷极限须分别定义：若对任意 $M>0$，都存在 $\delta>0$，使得当 $0<|x-x_0|<\delta$ 时有 $\alpha(x)>M$，则记为
\[
\lim_{x\to x_0}\alpha(x)=+\infty;
\]
若对任意 $M>0$，都存在 $\delta>0$，使得当 $0<|x-x_0|<\delta$ 时有 $\alpha(x)<-M$，则记为
\[
\lim_{x\to x_0}\alpha(x)=-\infty.
\]
上述两种情形均满足绝对值意义下的无穷大条件，但仅有 $|\alpha(x)|\to+\infty$ 不能推出其中任何一种带符号的极限。无穷小与无穷大都是相对于特定的极限过程而言的。\label{def:infinitesimal}
\end{definition}
% UNIT S03.U02
\begin{definition}
设函数 $\alpha(x)$ 与 $\beta(x)$ 在 $x_0$ 的某个去心邻域内有定义，且在该邻域内 $\beta(x)\neq0$。若
\[
\lim_{x\to x_0}\alpha(x)=0,\qquad
\lim_{x\to x_0}\beta(x)=0,
\qquad
\lim_{x\to x_0}\frac{\alpha(x)}{\beta(x)}=1,
\]
则称 $\alpha(x)$ 与 $\beta(x)$ 当 $x\to x_0$ 时为等价无穷小，记为
\[
\alpha(x)\sim\beta(x)\qquad(x\to x_0).
\]
其中，前两个条件说明二者都是同一极限过程下的无穷小，最后一个条件给出它们的等价关系；比值条件只在商有定义的去心邻域内使用。若只满足比值极限为 $1$ 而不要求二者都趋于零，则称为一般的等价函数关系，而不称为等价无穷小。\label{def:equivalent-infinitesimal}
\end{definition}
% UNIT S03.U03
\begin{theorem}
在实数范围内，当 $x\to0$ 时，有两个基本极限
\[
\lim_{x\to0}\frac{\sin x}{x}=1,
\qquad
\lim_{x\to0}\frac{e^x-1}{x}=1.
\]
其中，第一式中的 $x$ 取弧度制。由此可得
\[
\sin x\sim x,
\qquad
 e^x-1\sim x\qquad(x\to0).
\]
这些结论只适用于所写出的极限过程及其定义域条件，不能不加条件地推广到任意点或任意角度单位。\label{thm:fundamental-limits}
\end{theorem}
% UNIT S03.U04
\begin{proof}
先证明第一个基本极限。对 $0<x<\dfrac{\pi}{2}$，由单位圆中的几何不等式，有
\[
\sin x<x<\tan x.
\]
由于 $\sin x>0$，两边分别除以 $\sin x$，得到
\[
1<\frac{x}{\sin x}<\frac{1}{\cos x}.
\]
取倒数可得
\[
\cos x<\frac{\sin x}{x}<1.
\]
当 $x\to0^+$ 时，$\cos x\to1$，由夹逼定理可知
\[
\lim_{x\to0^+}\frac{\sin x}{x}=1.
\]
又因为 $\sin x/x$ 是偶函数，所以左极限也等于 $1$，从而
\[
\lim_{x\to0}\frac{\sin x}{x}=1.
\]

再证明第二个基本极限。这里使用指数函数的一个独立基本性质：对 $0<t<1$，有
\[
1+t\le e^t\le\frac{1}{1-t}.
\]
该不等式是指数函数基本性质的标准推论，独立于待证的差商极限，不由导数定义推出。

当 $0<x<1$ 时，由上述不等式得到
\[
1\le\frac{e^x-1}{x}\le\frac{1}{1-x}.
\]
令 $x\to0^+$，夹逼定理给出
\[
\lim_{x\to0^+}\frac{e^x-1}{x}=1.
\]
当 $x=-t<0$ 且 $0<t<1$ 时，对 $1+t\le e^t\le1/(1-t)$ 取倒数，得到
\[
1-t\le e^{-t}\le\frac{1}{1+t}.
\]
因此
\[
\frac{1}{1+t}\le\frac{e^{-t}-1}{-t}\le1.
\]
令 $t\to0^+$，再次由夹逼定理得到
\[
\lim_{x\to0^-}\frac{e^x-1}{x}=1.
\]
左右极限相等，所以
\[
\lim_{x\to0}\frac{e^x-1}{x}=1.
\]
整个论证只使用上述独立的指数函数不等式和夹逼定理，没有把待证极限作为指数函数在零点的导数值使用。
\end{proof}
% UNIT S03.U05
\begin{theorem}
设 $\alpha(x)$ 与 $\beta(x)$ 在 $x_0$ 的某个去心邻域内有定义，且在该邻域内均不为零。若
\[
\alpha(x)\sim\beta(x)\qquad(x\to x_0),
\]
并且函数 $F(x)$ 在相应去心邻域内有定义，满足 $F(x)\alpha(x)$ 有定义以及
\[
\lim_{x\to x_0}F(x)\alpha(x)=L,
\]
其中 $L$ 为有限数，则
\[
\lim_{x\to x_0}F(x)\beta(x)=L,
\]
只要 $F(x)\beta(x)$ 在该去心邻域内有定义。特别地，在乘积的极限中，可以用与某个无穷小等价的无穷小替换它，但必须保持相关表达式有定义。

商情形也有明确的替换条件：若 $F(x)/\alpha(x)$ 在某个去心邻域内有定义，且
\[
\lim_{x\to x_0}\frac{F(x)}{\alpha(x)}=L,
\]
其中 $L$ 为有限数，并且 $F(x)/\beta(x)$ 在同一去心邻域内有定义，则
\[
\lim_{x\to x_0}\frac{F(x)}{\beta(x)}=L.
\]
这里要求 $\alpha(x)$ 与 $\beta(x)$ 在该去心邻域内均不为零，从而 $\alpha(x)/\beta(x)$ 有定义；证明为
\[
\frac{F(x)}{\beta(x)}=\frac{F(x)}{\alpha(x)}\frac{\alpha(x)}{\beta(x)}.
\]
由 $\alpha(x)\sim\beta(x)$ 可知
\[
\lim_{x\to x_0}\frac{\alpha(x)}{\beta(x)}=1.
\]
再由极限的乘法法则，得到
\[
\lim_{x\to x_0}\frac{F(x)}{\beta(x)}
=\left(\lim_{x\to x_0}\frac{F(x)}{\alpha(x)}\right)
\left(\lim_{x\to x_0}\frac{\alpha(x)}{\beta(x)}\right)
=L\cdot1=L.
\]
因此，在商的极限中，可以将分母中的无穷小 $\alpha(x)$ 替换为与其等价的无穷小 $\beta(x)$，但必须同时检查两个商在同一去心邻域内有定义，并且原商的有限极限存在。一般而言，不能在加减式中无条件逐项替换等价无穷小；加减项的替换必须另行检验整体极限，不能仅由等价关系和表达式有定义推出。\label{thm:equivalent-substitution}
\end{theorem}
% UNIT S03.U06
\begin{proof}
由 $\alpha(x)\sim\beta(x)$ 的定义，有
\[
\lim_{x\to x_0}\frac{\alpha(x)}{\beta(x)}=1.
\]
又由等价无穷小替换定理的假设，$\alpha(x)$ 与 $\beta(x)$ 在同一去心邻域内均不为零，因此商 $\beta(x)/\alpha(x)$ 在该邻域内有定义，并且
\[
\frac{\beta(x)}{\alpha(x)}
=\left(\frac{\alpha(x)}{\beta(x)}\right)^{-1}.
\]
由于 $\lim_{x\to x_0}\alpha(x)/\beta(x)=1\neq0$，由倒数极限定理可得
\[
\lim_{x\to x_0}\frac{\beta(x)}{\alpha(x)}=1.
\]
在所给去心邻域内，利用乘积的代数分解，
\[
F(x)\beta(x)=F(x)\alpha(x)\frac{\beta(x)}{\alpha(x)}.
\]
由假设和极限的乘法法则，得到
\[
\lim_{x\to x_0}F(x)\beta(x)
=\left(\lim_{x\to x_0}F(x)\alpha(x)\right)
\left(\lim_{x\to x_0}\frac{\beta(x)}{\alpha(x)}\right)
=L\cdot1=L.
\]
因此，在上述条件下，若原式写成 $F(x)\alpha(x)$，可以把其中的因子 $\alpha(x)$ 替换为 $\beta(x)$，得到 $F(x)\beta(x)$。这一操作等价于将原式整体乘以修正比值 $\beta(x)/\alpha(x)$，而不是用该比值替换原式中的某个因子。替换时，$\alpha(x)$ 与 $\beta(x)$ 在同一去心邻域内均须非零，相关表达式须有定义，并且 $\lim_{x\to x_0}F(x)\alpha(x)$ 必须存在且为有限数。该结论适用于所证明的乘积结构，以及通过倒数改写后满足同样条件的商结构。它不表示可以在加减式中无条件逐项替换等价无穷小；例如加减项的替换必须另行检验整体极限，不能仅由等价关系和表达式有定义推出。
\end{proof}
% UNIT S03.U07
\begin{example}
计算
\[
\lim_{x\to0}\frac{(1+x)\sin x}{x}.
\]
要求说明等价无穷小替换的依据及适用条件。
\end{example}
% UNIT S03.U08
\par\noindent\textbf{解：}当 $x\to0$ 时，$\sin x\sim x$，并且 $1+x\to1$。在 $x\neq0$ 的去心邻域内，原式有定义，且
\[
\frac{(1+x)\sin x}{x}=(1+x)\frac{\sin x}{x}.
\]
由基本极限
\[
\lim_{x\to0}\frac{\sin x}{x}=1,
\qquad
\lim_{x\to0}(1+x)=1,
\]
可得
\[
\lim_{x\to0}\frac{(1+x)\sin x}{x}=1\cdot1=1.
\]
这里的替换只在 $x\to0$ 且 $x\neq0$ 的去心邻域内进行，未将原式在 $x=0$ 处补定义。\par
% UNIT S03.U09
\begin{exercise}
\label{mathagent:exercise:410f498083ff9a28d1a2acf9839e74554b434d3445561642fdfe5d3854f06971}
计算
\[
\lim_{x\to0}\frac{e^x-1}{x}\,(1-2x),
\]
并说明所使用的基本极限及乘法法则的条件。
\end{exercise}
\section{函数连续性与间断点}\label{sec:S04}
% UNIT S04.U01
\begin{definition}
设函数 $f$ 在点 $x_0$ 的某个邻域内有定义。若 $f$ 在 $x_0$ 处有定义，且
\[
\lim_{x\to x_0}f(x)=f(x_0)\text{，}
\]
则称函数 $f$ 在点 $x_0$ 处连续。等价地，$f$ 在 $x_0$ 处连续是指：对任意 $\varepsilon>0$，都存在 $\delta>0$，使得当 $|x-x_0|<\delta$ 时有
\[
|f(x)-f(x_0)|<\varepsilon.
\]
其中，条件 $f(x_0)$ 有定义不可省略。若 $f$ 在区间内每一点都连续，则称 $f$ 在该区间上连续。\label{def:point-continuity}
\end{definition}
% UNIT S04.U02
\begin{definition}
设 $f$ 在区间 $I$ 上有定义。若对每个内点 $x_0\in I$，函数 $f$ 在 $x_0$ 处连续；当 $I$ 含有端点时，函数在左端点处连续是指右极限等于函数值，在右端点处连续是指左极限等于函数值，则称 $f$ 在区间 $I$ 上连续。

特别地，当 $I=[a,b]$ 且 $a<b$ 时，$f$ 在 $[a,b]$ 上连续意味着：对每个 $x_0\in(a,b)$，有
\[
\lim_{x\to x_0}f(x)=f(x_0)\text{，}
\]
并且有
\[
\lim_{x\to a^+}f(x)=f(a)\text{，}
\qquad
\lim_{x\to b^-}f(x)=f(b)\text{。}
\]
\label{def:interval-continuity}
\end{definition}
% UNIT S04.U03
\begin{definition}
设函数 $f$ 在点 $x_0$ 的某个去心邻域内有定义。若 $f(x_0)$ 未定义，或 $f(x_0)$ 有定义但 $f$ 在 $x_0$ 处不连续，则称 $x_0$ 为 $f$ 的间断点。因此，间断点不要求属于函数的定义域。常见类型如下。

若极限 $\lim_{x\to x_0}f(x)$ 存在且为有限数，但 $f(x_0)$ 未定义或 $f(x_0)\neq\lim_{x\to x_0}f(x)$，则称 $x_0$ 为可去间断点。

若左、右极限都存在且有限，但不相等，即
\[
\lim_{x\to x_0^-}f(x)\neq\lim_{x\to x_0^+}f(x)\text{，}
\]
则称 $x_0$ 为跳跃间断点。

为说明无穷间断点，先定义单侧无穷极限。在上述去心邻域条件下，若对任意 $M>0$，都存在 $\delta>0$，使得当 $0<x-x_0<\delta$ 时有 $f(x)>M$，则记为
\[
\lim_{x\to x_0^+}f(x)=+\infty.
\]
若对任意 $M>0$，都存在 $\delta>0$，使得当 $0<x-x_0<\delta$ 时有 $f(x)<-M$，则记为
\[
\lim_{x\to x_0^+}f(x)=-\infty.
\]
若对任意 $M>0$，都存在 $\delta>0$，使得当 $0<x_0-x<\delta$ 时有 $f(x)>M$，则记为
\[
\lim_{x\to x_0^-}f(x)=+\infty.
\]
若对任意 $M>0$，都存在 $\delta>0$，使得当 $0<x_0-x<\delta$ 时有 $f(x)<-M$，则记为
\[
\lim_{x\to x_0^-}f(x)=-\infty.
\]
这里的 $\delta$ 均可取小到使相应单侧去心邻域包含在函数的定义域内。

若至少有一个单侧极限为 $+\infty$ 或 $-\infty$，即函数值沿相应方向趋于正无穷或负无穷，则称 $x_0$ 为无穷间断点。若函数在 $x_0$ 附近无界，但不存在趋于 $+\infty$ 或 $-\infty$ 的单侧极限，则不称为无穷间断点，而归入第二类间断点中的无界振荡情形。分类时必须先说明函数在相关去心邻域内有定义以及所考察的极限存在性条件。\label{def:discontinuity-types}
\end{definition}
% UNIT S04.U04
\begin{theorem}
设函数 $f$ 与 $g$ 在点 $x_0$ 处连续，且相关运算在 $x_0$ 处有定义。则函数
\[
f+g,
\qquad f-g,
\qquad fg
\]
在 $x_0$ 处连续。若 $g(x_0)\neq0$，并且在 $x_0$ 的某个邻域内 $g(x)\neq0$，则
\[
\frac{f}{g}
\]
在 $x_0$ 处连续。

若函数 $g$ 在 $x_0$ 处连续，函数 $f$ 在 $g(x_0)$ 处连续，并且复合函数 $f\circ g$ 在 $x_0$ 的某个邻域内有定义，则 $f\circ g$ 在 $x_0$ 处连续。\label{thm:continuity-algebra}
\end{theorem}
% UNIT S04.U05
\begin{proof}
由连续性的定义，函数 $f$ 与 $g$ 在 $x_0$ 处连续分别意味着
\[
\lim_{x\to x_0}f(x)=f(x_0)\text{，}
\qquad
\lim_{x\to x_0}g(x)=g(x_0)\text{。}
\]
根据极限的和、差与积法则，有
\[
\lim_{x\to x_0}(f(x)+g(x))=f(x_0)+g(x_0)=(f+g)(x_0)\text{，}
\]
\[
\lim_{x\to x_0}(f(x)-g(x))=f(x_0)-g(x_0)=(f-g)(x_0)\text{，}
\]
以及
\[
\lim_{x\to x_0}f(x)g(x)=f(x_0)g(x_0)=(fg)(x_0)\text{。}
\]
因此和、差与积在 $x_0$ 处连续。

若 $g(x_0)\neq0$ 且 $g$ 在 $x_0$ 的某个邻域内不为零，则由商的极限法则，
\[
\lim_{x\to x_0}\frac{f(x)}{g(x)}
=\frac{f(x_0)}{g(x_0)}
=\left(\frac{f}{g}\right)(x_0)\text{，}
\]
所以商在 $x_0$ 处连续。

最后设 $g$ 在 $x_0$ 处连续，$f$ 在 $g(x_0)$ 处连续，且复合函数在相应邻域内有定义。令 $\varepsilon>0$。由于 $f$ 在 $g(x_0)$ 处连续，存在 $\eta>0$，使得当 $|t-g(x_0)|<\eta$ 时，
\[
|f(t)-f(g(x_0))|<\varepsilon.
\]
由于 $g$ 在 $x_0$ 处连续，存在 $\delta>0$，使得当 $|x-x_0|<\delta$ 时，
\[
|g(x)-g(x_0)|<\eta.
\]
对这样的 $x$，取 $t=g(x)$；由复合函数在该邻域内有定义，且 $|t-g(x_0)|<\eta$，从而
\[
|f(g(x))-f(g(x_0))|<\varepsilon.
\]
又因为 $(f\circ g)(x_0)=f(g(x_0))$，所以 $f\circ g$ 在 $x_0$ 处连续。
\end{proof}
% UNIT S04.U06
\begin{theorem}
设函数 $g$ 在点 $x_0$ 处连续，函数 $f$ 在点 $g(x_0)$ 处连续，并且存在 $x_0$ 的某个邻域 $U$，使得对每个 $x\in U$，$g(x)$ 都属于 $f$ 的定义域。则复合函数
\[
(f\circ g)(x)=f(g(x))
\]
在点 $x_0$ 处连续。\label{thm:composition-continuity}
\end{theorem}
% UNIT S04.U07
\begin{proof}
令 $\varepsilon>0$。由于 $f$ 在 $g(x_0)$ 处连续，存在 $\eta>0$，使得当 $|t-g(x_0)|<\eta$ 时，
\[
|f(t)-f(g(x_0))|<\varepsilon.
\]
由于 $g$ 在 $x_0$ 处连续，存在 $\delta>0$，使得当 $|x-x_0|<\delta$ 时，
\[
|g(x)-g(x_0)|<\eta.
\]
根据假设，当 $x$ 属于该邻域 $U$ 时，$g(x)$ 属于 $f$ 的定义域；缩小 $\delta$ 后还可令 $|x-x_0|<\delta$ 的点属于 $U$。于是对满足 $|x-x_0|<\delta$ 的 $x$，令 $t=g(x)$，便有
\[
|t-g(x_0)|=|g(x)-g(x_0)|<\eta
\]
以及
\[
|f(g(x))-f(g(x_0))|<\varepsilon.
\]
而
\[
(f\circ g)(x_0)=f(g(x_0))\text{，}
\]
故
\[
|(f\circ g)(x)-(f\circ g)(x_0)|<\varepsilon.
\]
由点连续的定义，$f\circ g$ 在 $x_0$ 处连续。
\end{proof}
% UNIT S04.U08
\begin{example}
判断函数
\[
f(x)=\frac{x^2-1}{x-1}
\]
在 $x=1$ 处的连续性，并指出该点的间断类型。
\end{example}
% UNIT S04.U09
\par\noindent\textbf{解：}函数 $f$ 的定义域为 $\mathbb{R}\setminus\{1\}$，因此 $f(1)$ 未定义。当 $x\neq1$ 时，
\[
f(x)=\frac{(x-1)(x+1)}{x-1}=x+1\text{。}
\]
所以
\[
\lim_{x\to1}f(x)=\lim_{x\to1}(x+1)=2\text{。}
\]
极限存在且为有限数，但函数在 $x=1$ 处没有定义。因此 $x=1$ 是可去间断点。若补定义 $f(1)=2$，所得函数将在 $x=1$ 处连续。\par
% UNIT S04.U10
\begin{example}
设
\[
f(x)=
\begin{cases}
1+x,&x<0,\\
2+x^2,&x\geq0\text{。}
\end{cases}
\]
判断 $f$ 在 $x=0$ 处的连续性及间断类型。
\end{example}
% UNIT S04.U11
\par\noindent\textbf{解：}由左侧表达式可得
\[
\lim_{x\to0^-}f(x)=\lim_{x\to0^-}(1+x)=1\text{。}
\]
由右侧表达式可得
\[
\lim_{x\to0^+}f(x)=\lim_{x\to0^+}(2+x^2)=2\text{。}
\]
又因为 $f(0)=2$，所以右极限等于函数值，但左极限与右极限不相等。根据双侧极限存在的充要条件，$\lim_{x\to0}f(x)$ 不存在。因此 $f$ 在 $x=0$ 处不连续，且 $x=0$ 是跳跃间断点。\par
% UNIT S04.U12
\begin{exercise}
\label{mathagent:exercise:89dc01d13bdd47d62d3f0c96d9d7bd4201cee6446e15a7269f32977f2b24d58c}
设
\[
f(x)=
\begin{cases}
\dfrac{x^2-4}{x-2},&x\neq2,\\
a,&x=2\text{。}
\end{cases}
\]
求常数 $a$ 的值，使函数 $f$ 在 $x=2$ 处连续，并说明若 $a$ 取其他值时该点的间断类型。
\end{exercise}
\section*{练习答案}
% UNIT S01.U13
\par\noindent\textbf{\hyperref[mathagent:exercise:bf60311c135cf0c14f3cb42069a931ed175758be7d89dff861bed83a8d332803]{练习 \ref*{mathagent:exercise:bf60311c135cf0c14f3cb42069a931ed175758be7d89dff861bed83a8d332803} 的解答：}}当 $x\to1^-$ 时，$f(x)=2x+1$，因此
\[
\lim_{x\to1^-}f(x)=2\cdot1+1=3\text{。}
\]
当 $x\to1^+$ 时，$f(x)=5$，因此
\[
\lim_{x\to1^+}f(x)=5\text{。}
\]
由于左极限 $3$ 与右极限 $5$ 不相等，根据定理 \ref{thm:two-sided-limit-criterion}，双侧极限 $\lim_{x\to1}f(x)$ 不存在。\par
% UNIT S02.U13
\par\noindent\textbf{\hyperref[mathagent:exercise:b4c9c9a1fb82bcef052fdbf6e1539831ed88cee016dbcbf83a50e5e8d2311be0]{练习 \ref*{mathagent:exercise:b4c9c9a1fb82bcef052fdbf6e1539831ed88cee016dbcbf83a50e5e8d2311be0} 的解答：}}由极限的和与差法则，
\[
\lim_{x\to x_0}(f(x)+g(x))=2+(-1)=1,
\]
\[
\lim_{x\to x_0}(f(x)-g(x))=2-(-1)=3.
\]
分母的极限为 $3\neq0$，所以分母在 $x_0$ 的某个去心邻域内不为零，可以使用商的极限法则。因此
\[
\lim_{x\to x_0}\frac{f(x)+g(x)}{f(x)-g(x)}=\frac{1}{3}.
\]\par
% UNIT S03.U10
\par\noindent\textbf{\hyperref[mathagent:exercise:410f498083ff9a28d1a2acf9839e74554b434d3445561642fdfe5d3854f06971]{练习 \ref*{mathagent:exercise:410f498083ff9a28d1a2acf9839e74554b434d3445561642fdfe5d3854f06971} 的解答：}}当 $x\to0$ 且 $x\neq0$ 时，$\dfrac{e^x-1}{x}$ 有定义。由基本极限
\[
\lim_{x\to0}\frac{e^x-1}{x}=1,
\]
以及
\[
\lim_{x\to0}(1-2x)=1,
\]
再由极限的乘法法则，得到
\[
\lim_{x\to0}\frac{e^x-1}{x}(1-2x)=1\cdot1=1.
\]
乘法法则要求两个因子在同一去心邻域内有定义，并且各自的有限极限存在；这些条件在本题中均满足。\par
% UNIT S04.U13
\par\noindent\textbf{\hyperref[mathagent:exercise:89dc01d13bdd47d62d3f0c96d9d7bd4201cee6446e15a7269f32977f2b24d58c]{练习 \ref*{mathagent:exercise:89dc01d13bdd47d62d3f0c96d9d7bd4201cee6446e15a7269f32977f2b24d58c} 的解答：}}当 $x\neq2$ 时，
\[
\frac{x^2-4}{x-2}=\frac{(x-2)(x+2)}{x-2}=x+2\text{。}
\]
因此
\[
\lim_{x\to2}f(x)=\lim_{x\to2}(x+2)=4\text{。}
\]
函数在 $x=2$ 处连续的条件是
\[
\lim_{x\to2}f(x)=f(2)=a\text{，}
\]
所以必须取
\[
a=4\text{。}
\]
当 $a=4$ 时，函数在 $x=2$ 处连续；当 $a\neq4$ 时，极限仍存在且等于 $4$，但函数值不等于极限值，因此 $x=2$ 是可去间断点。\par
\end{document}
